%======================================================================
% Appendix C: Benchmark specifications
%======================================================================
\section{Benchmark Specifications}
\label{app:benchmarks}

The examples below define reproducible comparison problems and the quantities to specify. They do not state numerical outcomes.

\subsection{BM-1: Classical elastic interface}

\textbf{Problem:} Plane P-wave incident on a perfectly bonded interface between two elastic half-spaces.

\textbf{Governing equations:} Navier equations for the stated isotropic elastic model in each medium. Specify densities, Lam\'e moduli, incidence angle, normal orientation, and amplitude normalization before computing coefficients.

\textbf{Comparison:} Evaluate the complete time-averaged normal flux and include all propagating channels and relevant cross terms. The shorthand $R_P+T_P=1$ applies only when its single-channel and flux-normalization assumptions hold.

\subsection{BM-2: Classical thermoelastic interface}

\textbf{Problem:} Coupled thermoelastic wave at a welded interface.

\textbf{Comparison:} Specify the coupled constitutive model, boundary data, material parameters, and flux convention before comparing reflected and transmitted fields.

\subsection{BM-3: Piezothermoelastic interface}

\textbf{Problem:} Coupled elastic, thermal, and quasistatic electrical fields at an interface.

\textbf{Comparison:} State the constitutive equations and account for mechanical, electrical, and thermal power channels consistently.

\subsection{BM-4: One-dimensional nonlocal elastic dispersion}

\textbf{Problem:} For the specific Helmholtz-type model stated in Section~\ref{sec:nonlocal}, compare the dispersion relation and its long- and short-wavelength limits with the classical bar.

\subsection{BM-5: Memory-dependent thermoelastic wave}

\textbf{Problem:} Specify a memory kernel and initial-boundary-value problem; compare a numerical solution with an independently derived reference solution.
